% This is text associated with the open textbook "Open Electromagnetics"
% See immediately below for author, date
% This work is licensed under the Creative Commons Attribution-ShareAlike 4.0 International License. (CC BY-SA 4.0) 
% To view a copy of the license, visit http://creativecommons.org/licenses/by-sa/4.0/

%ORB TITLE Electric Field Due to Point Charges
%ORB AUTHOR 0        # S.W. Ellingson, Virginia Tech (ellingson@vt.edu)
%ORB DATE 20180803   # yyyymmdd
%ORB ABSTRACT Electric Field Due to Point Charges

\label{m0103_E_due_to_Point_Charges} % <-- Must match name of this file and module directory name.  Don't mess with this!

The electric field intensity associated with a single particle bearing charge $q_1$, located at the origin, is (Section~\ref{m0102_Coulombs_Law})
%
\begin{equation}
{\bf E}({\bf r}) = \hat{\bf r}\frac{q_1}{4\pi\epsilon r^2}
\end{equation} 
%
If this particle is instead located at some position ${\bf r}_1$, then the above expression may be written as follows:
%
\begin{equation}
{\bf E}({\bf r};{\bf r}_1) = \frac{{\bf r}-{\bf r}_1}{\left|{\bf r}-{\bf r}_1\right|}~\frac{q_1}{4\pi\epsilon \left|{\bf r}-{\bf r}_1\right|^2}
\end{equation} 
%
or, combining like terms in the denominator:
%
\begin{equation}
{\bf E}({\bf r};{\bf r}_1) = \frac{{\bf r}-{\bf r}_1}{\left|{\bf r}-{\bf r}_1\right|^3}~\frac{q_1}{4\pi\epsilon}
\end{equation} 

Now let us consider the field due to multiple such particles.
Under the usual assumptions about the permittivity of the medium (reminder: Section~\ref{m0007_Properties_of_Materials}), the property of superposition applies.
Using this principle, we conclude: 
%
%ORB HIGHLIGHT
\begin{mdframed}[backgroundcolor=yellow!20,nobreak=true] 
The electric field resulting from a set of charged particles is equal to the sum of the fields associated with the individual particles.
\end{mdframed}
%
Stated mathematically:
%
\begin{equation}
{\bf E}({\bf r}) = \sum_{n=1}^{N}{\bf E}({\bf r};{\bf r}_n)
\end{equation} 
%
where $N$ is the number of particles.
Thus, we have
%
\begin{equation}
{\bf E}({\bf r}) = \frac{1}{4\pi\epsilon} \sum_{n=1}^{N} {   \frac{{\bf r}-{\bf r}_n}{\left|{\bf r}-{\bf r}_n\right|^3}~q_n}
\end{equation} 

% TODO: complete example below.
% %ORB EXAMPLE
% ~\\ \begin{example} 
% \begin{mdframed}[backgroundcolor=green!20] 
% Electric Dipole
% \end{mdframed}
% {\color{red}NEEDS TO BE COMPLETED.}
% \centerline{\rule{2in}{1pt}} \end{example}





